% slide-ready summary for the lecture box 'Research from our group' (max 6 lines); generated by the notebook
\begin{itemize}\setlength{\itemsep}{2pt}
  \item \textbf{Criza din vara anului 2024} (Andrei, Pele și Cramer, 2026, PICBE): în scenariile contrafactuale estimate cu Random Forest, răspunsul cererii este pe primul loc; nicio intervenție nu coboară vîrful sub 200 EUR/MWh (p.~95).
  \item Datele cursului: criza s-a manifestat seara --- 916,72 EUR/MWh pe 18 iulie 2024, la ora 21:00; producția solară medie la ora 21:00 a fost de 1 MW.
  \item Replicarea noastră: o reducere a cererii cu 10\,\% scade prețul în orele de vîrf cu 16,9 EUR/MWh (RF), respectiv cu 16,2 EUR/MWh (Ridge); o producție solară mai mare cu 50\,\% nu schimbă aproape nimic (0,5 EUR/MWh).
  \item Concluzie: antrenat pe luni obișnuite, modelul RF subestimează prețurile de peste 500 EUR/MWh cu 59\,\% chiar și cu previziune perfectă; numai cu informația disponibilă la închiderea porții, MAE crește de la 32,4 la 35,4 EUR/MWh.
\end{itemize}
